\chapter{Analiza problemu i przegląd istniejących rozwiązań}
\label{cha:analiza_problemu}

Utrzymanie samochodu to nie tylko regularne tankowanie czy ładowanie, ale też serwis, ubezpieczenia, przeglądy i~raty leasingowe. Wielu kierowców zapisuje te koszty na papierowych paragonach, w~luźnych notatkach albo próbuje je po prostu zapamiętać. Takie rozproszenie utrudnia policzenie realnych wydatków na auto i~wyłapanie nagłego wzrostu spalania, który często sygnalizuje usterkę. Problem rośnie, gdy w~domu jest kilka pojazdów, a~o~kończącej się polisie OC czy badaniu technicznym łatwo zapomnieć.

% ==============================================================================
\section{Uzasadnienie potrzeby opracowania rozwiązania}
\label{sec:uzasadnienie_rozwiazania}

Dostępne na rynku aplikacje do kontroli wydatków na auto nie zawsze sprawdzają się w~codziennym użyciu. Wiele z~nich wymaga założenia konta i~stałego dostępu do Internetu, co utrudnia szybki zapis danych na stacji paliw ze słabym zasięgiem albo w~garażu podziemnym. Część darmowych narzędzi blokuje obsługę kilku pojazdów i~skanowanie paragonów za płatną subskrypcją lub zasypuje ekran reklamami.

Projektowana aplikacja \textbf{MotoLog} rozwiązuje te problemy, przenosząc zapis i~przetwarzanie danych bezpośrednio na telefon użytkownika. Dzięki temu system działa szybko, bez konieczności rejestracji i~niezależnie od połączenia z~siecią.

\subsection{Główne cele projektowanego systemu}
\label{subsec:glowne_cele}

Głównym celem jest uproszczenie zapisu tankowań i~ładowań dzięki odczytowi paragonów (OCR) oraz automatycznemu wyliczaniu średniego zużycia paliwa lub energii i~zasięgu auta. Aplikacja ma też porządkować wydatki poprzez podział na kategorie, historię wpisów oraz prognozę miesięcznego budżetu. Ważnym założeniem jest obsługa kilku pojazdów o~różnych napędach oraz wysyłanie lokalnych powiadomień, które chronią przed przegapieniem terminów opłat i~przeglądów.

Szczegółową analizę problemów biznesowych, zakres projektu oraz pełną listę wymagań i~reguł działania systemu opisano w~Rozdziale~\ref{cha:analiza_wymagan}.


% % ==============================================================================
% \section{Przegląd istniejących rozwiązań}
% \label{sec:przeglad_rozwiazan}

% \subsection{Analiza wybranych aplikacji}
% \label{subsec:analiza_aplikacji}

% \subsubsection{Rozwiązanie 1 (np. https://spalanie.formuly.pl/)}
% \label{subsubsec:rozwiazanie_1}
% % TODO: Opis, funkcjonalności, zalety oraz ograniczenia pierwszego rozwiązania.

% \subsubsection{Rozwiązanie 2 (np. Drivvo)}
% \label{subsubsec:rozwiazanie_2}
% % TODO: Opis, funkcjonalności, zalety oraz ograniczenia drugiego rozwiązania.

% \subsubsection{Rozwiązanie 3 (np. Yanosik)}
% \label{subsubsec:rozwiazanie_3}
% % TODO: Opis, funkcjonalności, zalety oraz ograniczenia trzeciego rozwiązania.


% \subsection{Porównanie analizowanych systemów}
% \label{subsec:porownanie_rozwiazan}

% Zestawienie porównawcze przeanalizowanych rozwiązań przedstawiono w~tabeli~\ref{tab:porownanie_rozwiazan}.

% \begingroup
% \renewcommand{\arraystretch}{1.3}
% \setlength{\LTcapwidth}{\linewidth}
% \begin{longtable}{|>{\centering\arraybackslash}m{2.8cm}|*{3}{m{\dimexpr(\linewidth - 2.8cm - 10\tabcolsep - 6\arrayrulewidth)/4\relax}|}>{\bfseries}m{\dimexpr(\linewidth - 2.8cm - 10\tabcolsep - 6\arrayrulewidth)/4\relax}|}
%     \caption{Porównanie istniejących rozwiązań z~projektowanym systemem}
%     \label{tab:porownanie_rozwiazan} \\
%     \hline
%     \rowcolor{tableheader}
%     \textbf{Cecha / Kryterium} & \multicolumn{1}{c|}{\textbf{Rozwiązanie A}} & \multicolumn{1}{c|}{\textbf{Rozwiązanie B}} & \multicolumn{1}{c|}{\textbf{Rozwiązanie C}} & \multicolumn{1}{c|}{\textbf{MotoLog (Projekt)}} \\
%     \hline
%     \endfirsthead

%     \multicolumn{5}{c}{\footnotesize\bfseries \tablename\ \thetable{} -- kontynuacja z~poprzedniej strony} \\
%     \hline
%     \rowcolor{tableheader}
%     \textbf{Cecha / Kryterium} & \multicolumn{1}{c|}{\textbf{Rozwiązanie A}} & \multicolumn{1}{c|}{\textbf{Rozwiązanie B}} & \multicolumn{1}{c|}{\textbf{Rozwiązanie C}} & \multicolumn{1}{c|}{\textbf{MotoLog (Projekt)}} \\
%     \hline
%     \endhead

%     \multicolumn{5}{r}{\footnotesize\emph{Ciąg dalszy na następnej stronie\ldots}} \\
%     \endfoot

%     \zrodloLT{5}{Opracowanie własne.}
%     \endlastfoot

%     \textbf{Główne funkcjonalności} & 
%     & 
%     & 
%     & \\
%     \hline
%     \textbf{Zalety} & 
%     & 
%     & 
%     & \\
%     \hline
%     \textbf{Ograniczenia} & 
%     & 
%     & 
%     & \\
%     \hline
%     \textbf{Działanie offline} & 
%     & 
%     & 
%     & \\
%     \hline
% \end{longtable}
% \endgroup

% % ==============================================================================
% \section{Wyróżniki projektowanego systemu}
% \label{sec:wyrozniki_systemu}

% % TODO: Elementy wyróżniające projektowany system na tle istniejących rozwiązań
% \begin{itemize}
%     \item \textbf{Element 1} -- % opis...
%     \item \textbf{Element 2} -- % opis...
%     \item \textbf{Element 3} -- % opis...
% \end{itemize}





% ==============================================================================
\section{Przegląd istniejących rozwiązań}
\label{sec:przeglad_rozwiazan}

W~niniejszym podrozdziale przeanalizowano pięć istniejących narzędzi wspierających kierowców w~kontroli zużycia paliwa oraz kosztów utrzymania pojazdu: kalkulator przeglądarkowy \textbf{spalanie.formuly.pl}, aplikacje mobilne \textbf{Drivvo} i~\textbf{Yanosik} oraz platformy webowo-mobilne \textbf{Fuelly} i~\textbf{Spritmonitor}.

\subsection{Analiza wybranych aplikacji}
\label{subsec:analiza_aplikacji}

% ------------------------------------------------------------------------------
\subsubsection{Rozwiązanie 1: spalanie.formuly.pl}
\label{subsubsec:rozwiazanie_1}

\textbf{Charakterystyka i~dostępność.} Serwis \textbf{spalanie.formuly.pl} to prosty kalkulator internetowy dostępny wyłącznie przez przeglądarkę WWW na komputerach i~urządzeniach mobilnych. Narzędzie nie posiada dedykowanych aplikacji natywnych w~sklepach Google Play i~App Store ani modułu rejestracji konta. Korzystanie ze strony jest bezpłatne, a~w~interfejsie mogą pojawiać się standardowe banery reklamowe.

\textbf{Zakres funkcji i~obsługa procesu.} Aplikacja działa w~modelu bezstanowym i~nie przechowuje historii wpisów ani profili pojazdów. Użytkownik wpisuje w~pojedynczym formularzu ilość zatankowanego paliwa, przejechany dystans oraz cenę za litr, a~system oblicza średnie spalanie na 100~km oraz koszt przejazdu. Serwis nie obsługuje wielu pojazdów, ewidencji wydatków serwisowych i~ubezpieczeniowych, wykresów trendów, przypomnień o~terminach, eksportu danych ani integracji z~innymi usługami.

\textbf{Ocena użyteczności i~prywatności.} Interfejs ogranicza się do jednego formularza, dzięki czemu obsługa jest natychmiastowa i~nie wymaga podawania żadnych danych osobowych. Brak lokalnej bazy wymusza jednak stały dostęp do Internetu przy każdym użyciu. Rozwiązanie sprawdza się wyłącznie jako doraźny kalkulator dla jednego przejazdu i~nie nadaje się do długoterminowej kontroli budżetu ani prowadzenia kilku aut.

\begin{figure}[H]
    \centering
    % \includegraphics[width=0.85\linewidth]{figures/spalanie_formuly_glowny.png}
    \caption{Widok główny kalkulatora w~serwisie spalanie.formuly.pl}
    \label{fig:spalanie_formuly_glowny}
    \zrodlo{Opracowanie własne na podstawie widoku aplikacji WWW (brak oddzielnych ekranów pojazdów, raportów, wydatków i~przypomnień ze względu na bezstanowy charakter narzędzia).}
\end{figure}

% ------------------------------------------------------------------------------
\subsubsection{Rozwiązanie 2: Drivvo}
\label{subsubsec:rozwiazanie_2}

\textbf{Charakterystyka i~dostępność.} \textbf{Drivvo} to rozbudowany dziennik eksploatacji pojazdu dostępny jako aplikacja mobilna na systemy Android i~iOS oraz poprzez panel przeglądarkowy WWW. Korzystanie z~systemu w~podstawowym zakresie jest darmowe po utworzeniu konta użytkownika (za pomocą adresu e-mail lub logowania społecznościowego), przy czym w~wersji bezpłatnej wyświetlane są reklamy. Producent udostępnia również płatny wariant rozszerzony, który usuwa reklamy oraz odblokowuje zaawansowane raportowanie, kopie zapasowe i~eksport danych.

\textbf{Zakres funkcji i~obsługa procesu.} Dodanie pojazdu wymaga określenia jego typu, marki, modelu, rodzaju paliwa oraz początkowego stanu licznika. System obsługuje wiele pojazdów w~ramach jednego profilu i~pozwala rejestrować tankowania z~automatycznym wyliczeniem średniego zużycia paliwa. Ewidencja kosztów obejmuje szeroką listę kategorii, w~tym naprawy serwisowe, wymianę oleju i~opon, ubezpieczenia, podatki, myjnię, parkingi oraz opłaty drogowe. Aplikacja generuje wykresy wydatków i~spalania w~czasie, wysyła przypomnienia powiązane z~datą lub przebiegiem oraz umożliwia eksport zestawień do plików CSV i~arkuszy kalkulacyjnych.

\textbf{Ocena użyteczności i~prywatności.} Bogaty zestaw funkcji analitycznych sprawia, że interfejs wymaga chwili na opanowanie struktury menu i~kategorii. Dane użytkowników są synchronizowane z~chmurą producenta, co ułatwia przeniesienie historii na nowe urządzenie, ale wiąże się z~przechowywaniem informacji poza telefonem. Dla zwykłego kierowcy wersja darmowa bywa uciążliwa z~powodu reklam, natomiast dla osób zarządzających kilkoma autami aplikacja oferuje bardzo szerokie możliwości ewidencyjne.

\begin{figure}[H]
    \centering
    % \includegraphics[width=0.45\linewidth]{figures/drivvo_glowny.png}
    % \includegraphics[width=0.45\linewidth]{figures/drivvo_pojazd.png}
    \caption{Aplikacja Drivvo: ekran główny oraz formularz dodawania pojazdu}
    \label{fig:drivvo_1}
    \zrodlo{Zrzuty ekranu z~aplikacji Drivvo.}
\end{figure}

\begin{figure}[H]
    \centering
    % \includegraphics[width=0.45\linewidth]{figures/drivvo_tankowanie.png}
    % \includegraphics[width=0.45\linewidth]{figures/drivvo_raporty.png}
    \caption{Aplikacja Drivvo: rejestracja tankowania oraz widok raportów i~statystyk}
    \label{fig:drivvo_2}
    \zrodlo{Zrzuty ekranu z~aplikacji Drivvo.}
\end{figure}

\begin{figure}[H]
    \centering
    % \includegraphics[width=0.45\linewidth]{figures/drivvo_wydatki.png}
    % \includegraphics[width=0.45\linewidth]{figures/drivvo_przypomnienia.png}
    \caption{Aplikacja Drivvo: ewidencja wydatków oraz konfiguracja przypomnień}
    \label{fig:drivvo_3}
    \zrodlo{Zrzuty ekranu z~aplikacji Drivvo.}
\end{figure}

% ------------------------------------------------------------------------------
\subsubsection{Rozwiązanie 3: Yanosik}
\label{subsubsec:rozwiazanie_3}

\textbf{Charakterystyka i~dostępność.} \textbf{Yanosik} to polska aplikacja mobilna na platformy Android oraz iOS, łącząca nawigację drogową i~system ostrzegania kierowców z~modułem obsługi pojazdu. Podstawowe funkcje drogowe są darmowe, jednak interfejs zawiera liczne banery reklamowe oraz oferty usług partnerskich. Dostępne są płatne pakiety rozszerzające funkcje asystenta jazdy i~ograniczające reklamy.

\textbf{Zakres funkcji i~obsługa procesu.} Korzystanie z~nawigacji nie wymaga logowania, natomiast zapisanie profilu pojazdu i~obsługa przypomnień wymagają założenia konta. W~sekcji pojazdu kierowca wprowadza numer rejestracyjny, markę, model oraz daty ważności polisy OC i~badania technicznego, z~możliwością dodania kilku samochodów. Aplikacja powiadamia o~zbliżających się terminach ubezpieczenia i~przeglądu oraz pozwala zapisywać podstawowe wydatki, jednak nie oferuje wyspecjalizowanej analityki spalania między pełnymi tankowaniami, skanowania paragonów OCR ani eksportu szczegółowych raportów kosztowych. System integruje się z~urządzeniami pokładowymi producenta oraz interfejsami Android Auto i~Apple CarPlay.

\textbf{Ocena użyteczności i~prywatności.} Ekran główny jest podporządkowany obsłudze w~trakcie jazdy, przez co funkcje typowo księgowe znajdują się głębiej w~menu. Aplikacja wymaga stałego połączenia z~siecią i~dostępu do lokalizacji, przetwarzając dane telemetryczne na serwerach operatora. Z~punktu widzenia zwykłego kierowcy jest to wygodny asystent uniwersalny, ale jako narzędzie do precyzyjnego rozliczania spalania i~kosztów wielu pojazdów wypada słabiej od dedykowanych dzienników.

\begin{figure}[H]
    \centering
    % \includegraphics[width=0.45\linewidth]{figures/yanosik_glowny.png}
    % \includegraphics[width=0.45\linewidth]{figures/yanosik_pojazd.png}
    \caption{Aplikacja Yanosik: ekran główny oraz moduł dodawania pojazdu}
    \label{fig:yanosik_1}
    \zrodlo{Zrzuty ekranu z~aplikacji Yanosik.}
\end{figure}

\begin{figure}[H]
    \centering
    % \includegraphics[width=0.45\linewidth]{figures/yanosik_tankowanie.png}
    % \includegraphics[width=0.45\linewidth]{figures/yanosik_raporty.png}
    \caption{Aplikacja Yanosik: rejestracja kosztów paliwa oraz podsumowanie wydatków}
    \label{fig:yanosik_2}
    \zrodlo{Zrzuty ekranu z~aplikacji Yanosik.}
\end{figure}

\begin{figure}[H]
    \centering
    % \includegraphics[width=0.45\linewidth]{figures/yanosik_wydatki.png}
    % \includegraphics[width=0.45\linewidth]{figures/yanosik_przypomnienia.png}
    \caption{Aplikacja Yanosik: lista wydatków oraz przypomnienia o~terminach OC i~przeglądu}
    \label{fig:yanosik_3}
    \zrodlo{Zrzuty ekranu z~aplikacji Yanosik.}
\end{figure}

% ------------------------------------------------------------------------------
\subsubsection{Rozwiązanie 4: Fuelly}
\label{subsubsec:rozwiazanie_4}

\textbf{Charakterystyka i~dostępność.} \textbf{Fuelly} to platforma do śledzenia zużycia paliwa i~kosztów serwisowych działająca poprzez serwis WWW oraz aplikacje mobilne na systemy Android i~iOS. Podstawowy zakres funkcji jest darmowy po rejestracji konta i~zawiera reklamy. Wariant płatny usuwa reklamy oraz rozszerza obsługę załączników fotograficznych, odczytu paragonów i~zaawansowanego raportowania.

\textbf{Zakres funkcji i~obsługa procesu.} Rejestracja wymaga utworzenia konta w~serwisie WWW lub aplikacji mobilnej. Profil pojazdu obejmuje rok produkcji, markę, model, wersję silnikową oraz wybrane jednostki miary, z~możliwością przypisania wielu aut do jednego konta. Dodawanie tankowania polega na wpisaniu przebiegu, ilości paliwa, ceny jednostkowej oraz oznaczeniu tankowania do pełna lub częściowego. Użytkownik może rejestrować wydatki eksploatacyjne i~serwisowe, analizować wykresy zużycia paliwa, ustawiać przypomnienia serwisowe oparte na dacie lub przebiegu oraz eksportować dane do plików CSV.

\textbf{Ocena użyteczności i~prywatności.} Interfejs kładzie nacisk na statystyki spalania i~porównywanie wyników z~innymi kierowcami, lecz barierą dla części polskich użytkowników jest anglojęzyczny charakter platformy. Wszystkie wpisy są synchronizowane z~centralną bazą portalu, co wymaga dostępu do sieci i~przekazywania danych na zewnętrzne serwery. System dobrze sprawdza się u~kierowców nastawionych na długoterminową analizę zużycia paliwa.

\begin{figure}[H]
    \centering
    % \includegraphics[width=0.45\linewidth]{figures/fuelly_glowny.png}
    % \includegraphics[width=0.45\linewidth]{figures/fuelly_pojazd.png}
    \caption{Platforma Fuelly: ekran główny oraz dodawanie pojazdu}
    \label{fig:fuelly_1}
    \zrodlo{Zrzuty ekranu z~aplikacji Fuelly.}
\end{figure}

\begin{figure}[H]
    \centering
    % \includegraphics[width=0.45\linewidth]{figures/fuelly_tankowanie.png}
    % \includegraphics[width=0.45\linewidth]{figures/fuelly_raporty.png}
    \caption{Platforma Fuelly: dodawanie tankowania oraz wykresy zużycia paliwa}
    \label{fig:fuelly_2}
    \zrodlo{Zrzuty ekranu z~aplikacji Fuelly.}
\end{figure}

\begin{figure}[H]
    \centering
    % \includegraphics[width=0.45\linewidth]{figures/fuelly_wydatki.png}
    % \includegraphics[width=0.45\linewidth]{figures/fuelly_przypomnienia.png}
    \caption{Platforma Fuelly: rejestr wydatków serwisowych oraz przypomnienia}
    \label{fig:fuelly_3}
    \zrodlo{Zrzuty ekranu z~aplikacji Fuelly.}
\end{figure}

% ------------------------------------------------------------------------------
\subsubsection{Rozwiązanie 5: Spritmonitor}
\label{subsubsec:rozwiazanie_5}

\textbf{Charakterystyka i~dostępność.} \textbf{Spritmonitor} to europejska platforma do ewidencji zużycia paliwa, energii elektrycznej oraz kosztów eksploatacji pojazdów, dostępna przez portal WWW oraz aplikacje mobilne na systemy Android i~iOS. Podstawowe funkcje są dostępne bezpłatnie po założeniu konta użytkownika i~wspierane reklamami, natomiast płatny pakiet rozszerzony usuwa reklamy oraz dodaje wybrane narzędzia raportowania i~obsługi dokumentów.

\textbf{Zakres funkcji i~obsługa procesu.} Po rejestracji konta użytkownik dodaje profil pojazdu, określając markę, model, typ silnika, rodzaj paliwa lub napędu elektrycznego oraz jednostki rozliczeniowe. Platforma obsługuje wiele pojazdów jednocześnie. Wpisy tankowań i~ładowań uwzględniają pełne i~częściowe uzupełnienia, a~moduł kosztów pozwala rejestrować ubezpieczenia, podatki, naprawy, wymianę opon oraz opłaty drogowe. System udostępnia szczegółowe statystyki spalania, porównanie wyników z~bazą innych użytkowników danego modelu, przypomnienia o~terminach oraz eksport danych do plików CSV.

\textbf{Ocena użyteczności i~prywatności.} Interfejs ma surowy, techniczny charakter, dzięki czemu wprowadzanie danych jest szybkie, choć mniej nowoczesne wizualnie. Dane są przechowywane na serwerach portalu i~mogą zasilać anonimowe zestawienia spalania dla konkretnych modeli aut. Rozwiązanie jest szczególnie wygodne dla kierowców prowadzących szczegółową ewidencję kilku pojazdów, w~tym aut elektrycznych i~hybrydowych.

\begin{figure}[H]
    \centering
    % \includegraphics[width=0.45\linewidth]{figures/spritmonitor_glowny.png}
    % \includegraphics[width=0.45\linewidth]{figures/spritmonitor_pojazd.png}
    \caption{System Spritmonitor: ekran główny oraz formularz dodawania pojazdu}
    \label{fig:spritmonitor_1}
    \zrodlo{Zrzuty ekranu z~aplikacji Spritmonitor.}
\end{figure}

\begin{figure}[H]
    \centering
    % \includegraphics[width=0.45\linewidth]{figures/spritmonitor_tankowanie.png}
    % \includegraphics[width=0.45\linewidth]{figures/spritmonitor_raporty.png}
    \caption{System Spritmonitor: dodawanie tankowania oraz statystyki spalania}
    \label{fig:spritmonitor_2}
    \zrodlo{Zrzuty ekranu z~aplikacji Spritmonitor.}
\end{figure}

\begin{figure}[H]
    \centering
    % \includegraphics[width=0.45\linewidth]{figures/spritmonitor_wydatki.png}
    % \includegraphics[width=0.45\linewidth]{figures/spritmonitor_przypomnienia.png}
    \caption{System Spritmonitor: ewidencja kosztów oraz przypomnienia o~terminach}
    \label{fig:spritmonitor_3}
    \zrodlo{Zrzuty ekranu z~aplikacji Spritmonitor.}
\end{figure}


% ------------------------------------------------------------------------------
\subsection{Porównanie analizowanych systemów}
\label{subsec:porownanie_rozwiazan}

Zestawienie porównawcze przeanalizowanych rozwiązań z~projektowanym systemem \textbf{MotoLog} przedstawiono w~tabeli~\ref{tab:porownanie_rozwiazan}.

\begingroup
\footnotesize
\renewcommand{\arraystretch}{1.3}
\setlength{\LTcapwidth}{\linewidth}
\begin{longtable}{|>{\centering\arraybackslash}m{2.3cm}|*{5}{m{\dimexpr(\linewidth - 2.3cm - 14\tabcolsep - 8\arrayrulewidth)/6\relax}|}>{\bfseries}m{\dimexpr(\linewidth - 2.3cm - 14\tabcolsep - 8\arrayrulewidth)/6\relax}|}
    \caption{Porównanie istniejących rozwiązań z~projektowanym systemem}
    \label{tab:porownanie_rozwiazan} \\
    \hline
    \rowcolor{tableheader}
    \textbf{Kryterium} & \multicolumn{1}{c|}{\textbf{spalanie.formuly.pl}} & \multicolumn{1}{c|}{\textbf{Drivvo}} & \multicolumn{1}{c|}{\textbf{Yanosik}} & \multicolumn{1}{c|}{\textbf{Fuelly}} & \multicolumn{1}{c|}{\textbf{Spritmonitor}} & \multicolumn{1}{c|}{\textbf{MotoLog (Projekt)}} \\
    \hline
    \endfirsthead

    \multicolumn{7}{c}{\footnotesize\bfseries \tablename\ \thetable{} -- kontynuacja z~poprzedniej strony} \\
    \hline
    \rowcolor{tableheader}
    \textbf{Kryterium} & \multicolumn{1}{c|}{\textbf{spalanie.formuly.pl}} & \multicolumn{1}{c|}{\textbf{Drivvo}} & \multicolumn{1}{c|}{\textbf{Yanosik}} & \multicolumn{1}{c|}{\textbf{Fuelly}} & \multicolumn{1}{c|}{\textbf{Spritmonitor}} & \multicolumn{1}{c|}{\textbf{MotoLog (Projekt)}} \\
    \hline
    \endhead

    \multicolumn{7}{r}{\footnotesize\emph{Ciąg dalszy na następnej stronie\ldots}} \\
    \endfoot

    \zrodloLT{7}{Opracowanie własne.}
    \endlastfoot

    \textbf{Dostępne platformy} & 
    Przeglądarka WWW & 
    Android, iOS, WWW & 
    Android, iOS & 
    Android, iOS, WWW & 
    Android, iOS, WWW & 
    Android (natywna aplikacja mobilna) \\
    \hline
    \textbf{Główne przeznaczenie} & 
    Jednorazowe obliczenie spalania & 
    Ewidencja wydatków i~serwisu auta & 
    Nawigacja i~ostrzeganie drogowe & 
    Dziennik tankowań i~analiza spalania & 
    Baza spalania, energii i~kosztów & 
    Lokalny dziennik tankowań, kosztów i~alertów \\
    \hline
    \textbf{Konto i~tryb pracy} & 
    Brak konta, wymaga sieci WWW & 
    Konto użytkownika, synchronizacja w~chmurze & 
    Konto dla usług pojazdu, wymaga sieci & 
    Konto użytkownika, synchronizacja w~chmurze & 
    Konto użytkownika, synchronizacja w~chmurze & 
    Brak konta, 100\% lokalna baza offline \\
    \hline
    \textbf{Obsługa wielu pojazdów} & 
    Brak zapisu pojazdów & 
    Wiele pojazdów w~profilu & 
    Wiele pojazdów w~koncie kierowcy & 
    Wiele pojazdów w~profilu & 
    Wiele pojazdów z~osobnymi statystykami & 
    Nielimitowane lokalne profile pojazdów \\
    \hline
    \textbf{Ewidencja spalania i~kosztów} & 
    Wyłącznie kalkulator spalania & 
    Pełne i~częściowe tankowania oraz szerokie kategorie kosztów & 
    Podstawowy zapis wydatków bez zaawansowanego spalania & 
    Szczegółowe spalanie i~wydatki serwisowe & 
    Spalanie (l i~kWh) oraz koszty eksploatacji & 
    Spalanie (l i~kWh), koszty z~kategoriami i~prognoza budżetu \\
    \hline
    \textbf{Skanowanie paragonów (OCR)} & 
    Brak & 
    Dostępne w~wybranych wariantach & 
    Brak dedykowanego skanera & 
    Dostępne w~wariancie płatnym & 
    Załączanie zdjęć dokumentów & 
    Wbudowany lokalny OCR z~korektą pól \\
    \hline
    \textbf{Przypomnienia o~terminach} & 
    Brak & 
    Powiadomienia według daty i~przebiegu & 
    Powiadomienia o~polisie OC i~przeglądzie & 
    Przypomnienia serwisowe po dacie lub przebiegu & 
    Przypomnienia o~opłatach i~przeglądach & 
    Lokalne powiadomienia z~domyślnym buforem 7~dni \\
    \hline
    \textbf{Eksport i~kopia danych} & 
    Brak & 
    Eksport CSV/arkusze oraz kopia w~chmurze & 
    Brak eksportu dziennika kosztów & 
    Eksport do plików CSV & 
    Eksport do plików CSV & 
    Lokalny zapis w~bazie urządzenia \\
    \hline
\end{longtable}
\endgroup

\subsection{Wskazanie najlepszych rozwiązań w~kluczowych kategoriach}
\label{subsec:najlepsze_rozwiazania_kategorie}

Na podstawie analizy funkcjonalnej zestawionych narzędzi wyłoniono rozwiązania najlepiej sprawdzające się w~poszczególnych obszarach zastosowań:

\begin{itemize}
    \item \textbf{Monitorowanie spalania:} \textbf{Spritmonitor} oraz \textbf{Fuelly} -- oba systemy precyzyjnie rozróżniają tankowania pełne i~częściowe, obsługują różne jednostki zużycia (w~tym energię elektryczną w~Spritmonitor) oraz pozwalają zestawiać własne wyniki ze średnią dla danego modelu pojazdu.
    \item \textbf{Zarządzanie kosztami pojazdu:} \textbf{Drivvo} -- oferuje najbardziej rozbudowany podział kategorii wydatków eksploatacyjnych, serwisowych i~administracyjnych wraz z~czytelnymi wykresami struktury kosztów.
    \item \textbf{Flota wielu pojazdów:} \textbf{Spritmonitor} oraz \textbf{Drivvo} -- umożliwiają prowadzenie oddzielnej ewidencji tankowań, kosztów i~harmonogramów serwisowych dla wielu pojazdów w~jednym koncie użytkownika.
    \item \textbf{Darmowe funkcje:} \textbf{Spritmonitor} (w~zakresie pełnego dziennika spalania, kosztów i~przypomnień dostępnych bezpłatnie po rejestracji) oraz \textbf{spalanie.formuly.pl} (jako najszybsze narzędzie do jednorazowego obliczenia bez zakładania konta i~instalacji aplikacji).
    \item \textbf{Stosunek funkcji do ceny:} \textbf{Spritmonitor} -- zapewnia kompletny zestaw narzędzi do analizy zużycia paliwa i~energii, obsługi wielu pojazdów, przypomnień oraz eksportu danych CSV już w~wariancie bezpłatnym, bez blokowania podstawowej ewidencji za opłatą.
\end{itemize}


% ==============================================================================
\section{Wyróżniki projektowanego systemu}
\label{sec:wyrozniki_systemu}

Na tle przeanalizowanych rozwiązań projektowany system \textbf{MotoLog} wyróżnia się następującymi cechami wynikającymi ze specyfikacji wymagań (Rozdział~\ref{cha:analiza_wymagan}):

\begin{itemize}
    \item \textbf{Pełna niezależność od sieci i~brak wymogu rejestracji (offline-first)} -- wszystkie operacje zapisu, odczytu i~obliczeń realizowane są w~100\% lokalnie w~bazie danych na urządzeniu mobilnym, bez zakładania konta i~przesyłania danych o~pojazdach na zewnętrzne serwery.
    \item \textbf{Lokalny moduł OCR paragonów dostępny bez opłat} -- aplikacja umożliwia automatyczne odczytywanie danych tankowania ze zdjęcia paragonu bezpośrednio na telefonie, z~możliwością wielokrotnego powtórzenia zdjęcia oraz ręcznej korekty pól przed zapisem.
    \item \textbf{Automatyczna walidacja spójności danych} -- system na bieżąco przelicza cenę jednostkową i~koszt całkowity, blokuje wpisanie przebiegu mniejszego od poprzedniego oraz ostrzega przed rejestracją daty wstecznej mogącej zaburzyć statystyki i~prognozy.
    \item \textbf{Obsługa wielu napędów i~proaktywne przypomnienia} -- jeden interfejs obsługuje pojazdy spalinowe, hybrydowe i~elektryczne (l/100~km oraz kWh/100~km), wylicza przewidywany zasięg, prognozuje miesięczny budżet oraz automatycznie tworzy lokalne powiadomienia na 7~dni przed terminem opłaty lub przeglądu.
\end{itemize}